\documentclass[11pt,oneside]{report}
% ======================================================================
%  THE TRIADIC MEASUREMENT SUBSTRATE --- COMPLETE CORPUS (Volumes 1 & 2)
%  Aaron Switzer, 2026
%  Single-file build. Compile: pdflatex x3 (third pass settles the TOC).
% ======================================================================
% ======================================================================
%  TMS Volume 1 -- shared preamble
%  Compiles with pdfLaTeX both locally and on Overleaf (no exotic fonts).
% ======================================================================
\usepackage[T1]{fontenc}
\usepackage[utf8]{inputenc}
\usepackage{amsmath}
\usepackage{amssymb}
\usepackage{mathptmx}            % Times text + math (widely available)
\usepackage[scaled=0.90]{helvet} % sans for headings
\usepackage{courier}

\usepackage[letterpaper,margin=1in]{geometry}
\usepackage{amsthm}
\usepackage{mathtools}
\usepackage{bm}
\usepackage{microtype}
\usepackage{enumitem}
\usepackage{booktabs}
\usepackage{array}
\usepackage{longtable}
\usepackage{caption}
\usepackage{xcolor}
\usepackage{titlesec}
\usepackage{fancyhdr}
\usepackage[hidelinks]{hyperref}
\usepackage{tocbibind}

% ---- spacing -----------------------------------------------------------
\setlength{\parindent}{1.2em}
\setlength{\parskip}{0.35em}
\linespread{1.04}
\setlist{leftmargin=1.6em,topsep=0.25em,itemsep=0.15em,parsep=0pt}

% ---- theorem environments ---------------------------------------------
\newtheorem{theorem}{Theorem}[section]
\newtheorem{lemma}[theorem]{Lemma}
\newtheorem{corollary}[theorem]{Corollary}
\newtheorem{proposition}[theorem]{Proposition}

\theoremstyle{definition}
\newtheorem{definition}[theorem]{Definition}
\newtheorem{axiom}{Axiom}
\newtheorem{claim}[theorem]{Claim}
\newtheorem{principle}[theorem]{Principle}
\newtheorem{conjecture}[theorem]{Conjecture}
\newtheorem{prediction}[theorem]{Prediction}
\newtheorem{property}[theorem]{Property}
\newtheorem{postulate}[theorem]{Postulate}

\theoremstyle{remark}
\newtheorem{remark}[theorem]{Remark}
\newtheorem*{remark*}{Remark}

% ---- section heading look (unnumbered; author supplies the numbers) ----
\titleformat{\section}[hang]{\normalfont\large\bfseries\sffamily}{}{0pt}{}
\titleformat{\subsection}[hang]{\normalfont\normalsize\bfseries\sffamily}{}{0pt}{}
\titleformat{\subsubsection}[hang]{\normalfont\normalsize\bfseries\itshape}{}{0pt}{}
\titlespacing*{\section}{0pt}{1.4em}{0.5em}
\titlespacing*{\subsection}{0pt}{1.0em}{0.35em}
\titlespacing*{\subsubsection}{0pt}{0.8em}{0.3em}

% chapter (= each paper) heading
\titleformat{\chapter}[display]
  {\normalfont\sffamily\bfseries}
  {}{0pt}{\Huge}
\titlespacing*{\chapter}{0pt}{10pt}{24pt}

% ---- page-break quality: avoid stranded headings and orphaned/widowed lines ----
% (applies to both volumes via the shared preamble)
\widowpenalty=10000        % no single line at the top of a page
\clubpenalty=10000         % no single line at the bottom of a page
\displaywidowpenalty=10000 % same, around displayed equations
\brokenpenalty=4000        % discourage page breaks right after a hyphenated line
% titlesec: forbid a page break immediately after a heading (heading stranded at page foot)
\usepackage{needspace}
\usepackage{etoolbox}
\pretocmd{\section}{\Needspace*{4\baselineskip}}{}{}
\pretocmd{\subsection}{\Needspace*{3\baselineskip}}{}{}

% ---- running heads -----------------------------------------------------
\pagestyle{fancy}
\fancyhf{}
\fancyhead[L]{\small\itshape The Triadic Measurement Substrate}
\fancyhead[R]{\small\itshape\nouppercase{\rightmark}}
\fancyfoot[C]{\small\thepage}
\renewcommand{\headrulewidth}{0.4pt}
\fancypagestyle{plain}{\fancyhf{}\fancyfoot[C]{\small\thepage}\renewcommand{\headrulewidth}{0pt}}

% ---- abstract / keywords / references list ----------------------------
\newenvironment{paperabstract}
  {\begin{center}\begin{minipage}{0.88\textwidth}\small
   \noindent\textbf{Abstract.}\itshape\space}
  {\end{minipage}\end{center}\vspace{0.4em}}

\newcommand{\keywords}[1]{\noindent{\small\textbf{Keywords:} #1}\vspace{0.4em}}

% per-paper APA reference list (hanging indent), kept as authored
\newenvironment{reflist}
  {\par\small\setlength{\parindent}{0pt}%
   \setlength{\leftskip}{1.6em}\setlength{\parindent}{-1.6em}%
   \setlength{\parskip}{0.4em}%
   \raggedright\hyphenpenalty=50\exhyphenpenalty=50\relax
   \sloppy}
  {\par}

% convenience
\providecommand{\tightlist}{\setlength{\itemsep}{0pt}\setlength{\parskip}{0pt}}
\providecommand{\TMS}{\textsc{tms}}
\hyphenation{con-fir-ma-tion sub-strate}
\begin{document}

\thispagestyle{empty}
\begin{titlepage}
\centering
\vspace*{2cm}
{\sffamily\bfseries\Large The Triadic Measurement Substrate\par}
\vspace{0.6cm}
{\large $\bar{B} = 3\bar{\alpha}$\par}
\vspace{1cm}
{\sffamily\large Volume 1\par}
\vspace{1.4cm}
{\Large \bfseries Paper 1: Emergent 3D Information Topology and the Binary Alphabet via [[4,2,2]] Quantum Error Detection Structural Correspondence\par}
\vfill
{\large Aaron Switzer\par}
\vspace{0.4cm}
{\large 2026\par}
\vspace{1.2cm}
{\small\itshape Excerpted from the complete two-volume corpus, \emph{The Triadic Measurement Substrate}. Full corpus and reproducibility package: \url{https://doi.org/10.5281/zenodo.22035422}\par}
\end{titlepage}
\cleardoublepage
\pagenumbering{arabic}

% ---------------- Volume 1 / paper1 ----------------
\chapter*{Paper 1}
\markboth{Paper 1}{Paper 1}
\addcontentsline{toc}{chapter}{Paper 1: Emergent 3D Information Topology and the Binary Alphabet via $[[4,2,2]]$ Quantum Error Detection Structural Correspondence}
\begingroup\centering
{\large\itshape Emergent 3D Information Topology and the Binary Alphabet via $[[4,2,2]]$ Quantum Error Detection Structural Correspondence\par}\vspace{0.8em}
{\normalsize Aaron Switzer\par}
{\small May 2026\par}
\endgroup\vspace{1.0em}

\noindent{\small\textbf{Keywords:} Information Theory, FCC Lattice, $[[4,2,2]]$ Code, Geometric Necessity, Causal Succession, Planck Length, Binary Alphabet, Generalized Uncertainty Principle, Computational Irreducibility}\par\vspace{0.6em}

\section*{Status of Results}
\addcontentsline{toc}{section}{Status of Results}

This paper distinguishes three categories of claim, each labeled where
it appears in the text.

\textbf{Derived.} Results that follow from the five axioms and the
Principle of Generative Economy by explicit construction or proof. The
triadic minimum \ensuremath{\bar{B}} = \ensuremath{3\bar{\alpha}_{2}} (\S\,2.1), the hexagonal packing of \ensuremath{M\alpha} (\S\,2.2), FCC
necessity over HCP via the causal distinguishability argument (\S\,2.4),
the \ensuremath{1\to 4} multiplier and its remainder-free integer division (\S\,4.2), the
[[4,2,2]] structural correspondence and code distance 2 (\S\,V),
the wave equation \ensuremath{\partial ^{2}\varphi /\partial t^{2}} \ensuremath{-} \ensuremath{c^{2}\nabla ^{2}\varphi} = 0 in the continuum limit (\S\,VI), and
the propagation speed \ensuremath{c_{\mathrm{TMS}}} = \ensuremath{1/\sqrt{2}} from FCC face-diagonal geometry (\S\,6.5)
are derived in this sense. The scale-invariant k = 3n stabilizer
null-space law (\S\,5.4) is derived here as a structural property of FCC +
[[4,2,2]] dynamics, empirically verified through Paper 5
Appendix E at levels 1 and 2. The sublattice bifurcation corollary
(\S\,2.4) is derived here in skeleton form, with the full geometric
development and empirical verification in Paper 5 \S\,6.1.3 and Appendix E.
The conjugate pair \ensuremath{(\varphi ,} \ensuremath{\Pi )} and the second-order dynamics are derived from
signal conservation under FCC connectivity (\S\,6.3), with the uniqueness
of the second equation tightened in that section.

\textbf{Structural correspondence.} Results in which a TMS construction
faithfully reproduces the skeleton of a known framework without claiming
reproduction of every feature of that framework. The [[4,2,2]]
mapping of \S\,V is the canonical example: the stabilizer algebra, the two
logical degrees of freedom, and the code distance are reproduced
exactly; the full quantum amplitude formalism is deferred to Paper 2 and
subsequent work. The identification \ensuremath{\ell_{0}} \ensuremath{\equiv} \ensuremath{\ell_{p}} in \S\,3.2 is a structural
correspondence in the same sense --- the minimum resolvable length of
the emergent geometry is identified with the Planck length, with the
derivation of \ensuremath{\hbar} and G from the confirmation field dynamics deferred
(\S\,3.2 Remark).

\textbf{Predictions / Conjectures.} Falsifiable claims that follow from
the framework but whose closed forms are not derived within this paper.
The GUP scaling \ensuremath{\beta} = f(d, \ensuremath{K_{\psi})} (\S\,7.2), the computational irreducibility
connection of \S\,7.2a, and the unified geometric account of decoherence
(\S\,7.3) are predictions in this sense. They are consequences of the
framework's geometric architecture; they are not derivations of f, of
explicit decoherence rates, or of falsification thresholds. The
functional form of f is refined across Papers 2--4 and pinned to
numerical content in Paper 5's simulation work.

This three-way labeling applies throughout the volume. Every claim in
the body text falls into exactly one category, and every category's
status is visible at the point of claim.

\section*{Notation}
\addcontentsline{toc}{section}{Notation}

Throughout this paper, the bar notation distinguishes confirmed (locked)
elements from latent (unlocked) elements. \ensuremath{\bar{\alpha}} denotes an agent
participating in a confirmation event; \ensuremath{\bar{B}} denotes a confirmed
(actualized) bit. Unbarred \ensuremath{\alpha} and B denote their latent counterparts.
Dimensional subscripts \ensuremath{\bar{\alpha}_{2}} and \ensuremath{\bar{\alpha}_{3}} distinguish 2D and 3D confirmation
contexts respectively. A confirmed bit carries a binary value \ensuremath{\bar{B}} \ensuremath{\in} \{0,
1\}; its signed representation \ensuremath{\chi} \ensuremath{\equiv} \ensuremath{2\bar{B}} \ensuremath{-} 1 \ensuremath{\in} \{\ensuremath{-1,} +1\} is referred to as
its polarization.

\begin{paperabstract}
We propose the Triadic Measurement Substrate (TMS), a generative model
in which wave dynamics and a binary computational alphabet emerge
necessarily from the interaction between a Maximum Entropy Stochastic
Field \ensuremath{(U_{\mathrm{sh}})} and a 2D hexagonal agent manifold \ensuremath{(M\alpha ).} Beginning from five
axioms and the Principle of Generative Economy, we demonstrate that
hexagonal close-packing, the \ensuremath{1\to 4} informational multiplier, FCC lattice
geometry, the binary alphabet \ensuremath{\bar{B}} \ensuremath{\in} \{0, 1\} of confirmed bit values, the
conjugate field structure, and the [[4,2,2]] quantum
error-detection structural correspondence follow by geometric necessity.
The fundamental agent spacing is identified with the Planck length \ensuremath{\ell_{p},}
the flop time with the Planck time \ensuremath{t_{p},} and the lattice propagation speed
\ensuremath{c_{\mathrm{TMS}}} = \ensuremath{1/\sqrt{2}} is forced by the geometry of FCC face-diagonal connectivity.
The continuum limit reproduces the standard wave equation with speed
c.~The [[4,2,2]] structural correspondence is shown to
faithfully encode both the value of every confirmed bit and the
inheritance of that value through the causal stack --- two independent
logical degrees of freedom jointly protected by the two stabilizer
generators against any single-agent disruption from \ensuremath{U_{\mathrm{sh}}.} The TMS
contains no free \emph{dimensionless} numerical parameters --- every
dimensionless quantity is fixed by the axioms and the Principle of
Generative Economy --- and takes a single dimensionful input (one action
scale, \ensuremath{\hbar}), from which the speed of light and Newton's constant follow
(\S\,3.2--3.3). The absolute length scale is fixed by identifying the
substrate's minimum length with the Planck length. The wave equation is
not fitted to the model --- it is the model's necessary large-scale
behavior. The paper closes with a novel falsifiable prediction: the GUP
deformation
parameter \ensuremath{\beta} is not a universal constant but scales with the causal depth
and polarization complexity of the confirmation record, providing a
geometric and computational derivation of both the Generalized
Uncertainty Principle and quantum decoherence from a single structural
property of the FCC causal stack.

This paper constitutes Volume 1, Paper 1 of a larger work. Subsequent
papers in Volume 1 develop the computational dynamics of the binary
alphabet, the emergence of learning and hierarchy, subsystem
development, and simulation. Volume 2 addresses the agent-side
complement.
\end{paperabstract}

\section*{I. Definitions and Ontological Primitives}
\addcontentsline{toc}{section}{I. Definitions and Ontological Primitives}

\subsection*{The Principle of Generative Economy}

All structure selects the minimum sufficient configuration that
maximizes generative potential. This principle governs all selection
throughout the TMS --- of primitives, confirmation geometry, and lattice
structure.

This principle belongs to the deep tradition in physics of
parsimony-as-law. Hamilton's Principle of Least Action (Hamilton, 1834)
holds that nature selects the path that extremizes the action integral.
Jaynes' Maximum Entropy Principle (Jaynes, 1957) holds that unbiased
inference selects the least-structured distribution consistent with
known constraints. Occam's Razor, in its formal instantiation as Minimum
Description Length (Rissanen, 1978), holds that the simplest model
consistent with the data is preferred. These principles share a common
structure: nature does not take the long way home.

The Principle of Generative Economy is a specific instantiation of this
tradition, distinguished by a dual constraint. Existing parsimony
principles minimize in one direction --- cost, assumptions, description
length. The Principle of Generative Economy jointly minimizes structural
complexity and maximizes generative capacity. This dual constraint is
necessary because pure minimization selects the degenerate case: a point
is simpler than a sphere, but generates nothing; a line is simpler than
a triangle, but closes nothing. The sphere and the triangle are not
selected because they are small --- they are selected because they are
the minimum configurations that remain generatively non-trivial.
Kauffman's adjacent possible (Kauffman, 1993) is the closest conceptual
relative: the idea that systems naturally explore the maximum frontier
of what their current structure makes accessible. A parallel framing of
dual parsimony has been developed independently by Ma, Tsao, and Shum
(2022) for the emergence of intelligence, where Parsimony and
Self-Consistency together govern what to learn and how to learn. The
Principle of Generative Economy is the geometric form that the parsimony
tradition takes when the only primitive is the sphere and the only
operation is confirmation. In giving condition~(i) below a specific
computable form --- the moduli dimension of a configuration under
axiom-forced constraints --- the framework instantiates this established
tradition rather than inventing a new measure of simplicity: the reading
of complexity as a count of free parameters, and of parsimony as its
minimization, is standard across the description-length and
model-selection literature cited above. What is particular to this work
is not the measure but its application, and specifically the requirement
that the constraint set be derived from the axioms, developed below.

\textbf{Formal Statement (Principle of Generative Economy).} Let K
denote the set of constraints imposed on a given selection by the axioms
and by prior structure relevant to that selection. Let \ensuremath{\mathcal{C}_{K}} denote the
set of configurations satisfying K. A configuration C* \ensuremath{\in} \ensuremath{\mathcal{C}_{K}} is the
PGE-selection if and only if:

\emph{(i) Minimum specification.} C* minimizes the \emph{specification cost}
over \ensuremath{\mathcal{C}_{K}}, where the specification cost of a configuration is the number of
free real parameters required to fix it beyond what K forces --- that is, the
dimension of the space of K-satisfying configurations equivalent to it under the
symmetry K permits (its moduli dimension). For discrete selections this reduces
to an integer count of independent structural commitments; for continuous
selections it is the dimension of the orientation-and-shape moduli. The measure
is standard and computable: for example, in the selection of the geometric
primitive the sphere has specification cost zero (isotropic: no orientation or
shape parameter), while a cube costs three, an ellipsoid five, and an irregular
form more --- so the sphere is the unique minimizer.

\emph{(ii) Generative non-triviality.} C* admits downstream structure
under the substrate dynamics. It is not a generative dead-end.

When K admits a unique configuration satisfying both (i) and (ii), that
configuration is the unique PGE-selection. When (i) admits multiple
candidates, (ii) is the tiebreaker. When the unique
minimum-specification candidate fails (ii), K is too restrictive for
generative non-triviality and the axioms must supply further constraint.

\textbf{The constraint-set requirement.} The entire evidential content of a
PGE-selection lies not in the parameter count, which is unambiguous, but in the
specification of K. Accordingly this work adopts a strict requirement: at every
invocation, K must be \emph{derived from the axioms and declared before the
selection}, never reverse-engineered from the desired configuration. A
PGE-selection is rigorous exactly to the extent that its K is exhibited as
axiom-forced. As a worked demonstration on a contested case, the selection of
FCC over hexagonal close-packing (\S\,2.4) uses K = \{Axiom~4 (no overwrite)
together with the definition of a confirmation datum as a state individuated by
its causal relations\}; the causal-distinguishability constraint that forbids the
ABAB stacking is not an added assumption but a consequence of Axiom~4, since two
causally indistinguishable layers are, by the definition of the record, the same
recorded state --- which Axiom~4 forbids. Establishing the axiom-forced character
of K at each remaining invocation, grouped by constraint-type, is the principal
open formal task of the program; the selection functional itself is fixed and
computable.

\begin{remark}[On the Two Conditions]
Condition (i) is the
substantive content of \emph{minimum sufficient}: sufficiency is set by
K, minimality is the moduli dimension not forced by K. Condition (ii) is
the substantive content of \emph{maximizes generative potential}:
generatively sterile candidates are excluded even when minimal. The pair
is necessary because pure minimization selects the degenerate case as
noted above. Every subsequent invocation of ``by Generative Economy'' in
this volume is to be verified against this formal statement --- the
constraint set K must be specifiable and axiom-derived, and C* must
satisfy both (i) and (ii) uniquely within \ensuremath{\mathcal{C}_{K}.}
\end{remark}

\textbf{Illustration (The Sphere Primitive).} As a first application:
for the selection of the geometric primitive in \S\,1.2, K = \{isotropy,
geometric primitive, supports the measurement role required by Axiom
1\}. Within \ensuremath{\mathcal{C}_{K},} the sphere has zero orientation parameters; any other
candidate (cube, ellipsoid, regular polyhedron, irregular form)
introduces at least one orientation or shape parameter not forced by K.
The sphere therefore satisfies (i) uniquely. It satisfies (ii) --- it
encloses interstices, supports contact geometry, packs in 2D and 3D, and
admits the cascade dynamics of subsequent sections. The sphere is the
unique PGE-selection. Axiom 2 records this result.

\subsection*{Axiom 0 --- Information Fundamentalism}

Information is fundamental. It exists as potential prior to
confirmation. This potential field is \ensuremath{U_{\mathrm{sh}}} --- pure informational
superposition, algorithmically incompressible, with complexity K(n) \ensuremath{\approx} n
(Kolmogorov, 1965). It is maximum entropy by definition: a unique state,
not a tunable quantity. \ensuremath{U_{\mathrm{sh}}} has no free parameters. This axiom is the
formal expression of the informational priority articulated by Wheeler's
``it from bit'' (Wheeler, 1990): the physical arises from the
informational, with binary distinctions as the deepest layer of
substrate from which all structure emerges.

\subsection*{Axiom 1 --- The Measurement Imperative}

Information must be measured to have meaning. Therefore a measurement
substrate must exist.

\subsection*{Axiom 2 --- The Sphere as Primitive}

Given the Principle of Generative Economy, the unique minimum sufficient
measurement substrate is the sphere --- the only isotropic primitive
requiring zero orientation parameters. The circle is its 2D expression.
Hexagonal close-packing and the FCC lattice follow as the unique minimum
sufficient tilings of this primitive in 2D and 3D respectively.

\subsection*{Axiom 3 --- Signal Conservation}

Informational signals are conserved across flops. They cannot be
created, destroyed, or deferred.

\subsection*{Axiom 4 --- Causal Anchoring}

Confirmed states cannot be overwritten. A confirmed state is the
geometric record of the signals that produced it --- to overwrite it
would violate signal conservation. The system therefore has only one
available direction: forward. Each flop adds a new layer to the causal
record without disturbing prior layers. This directional necessity is
the geometric origin of flop time \ensuremath{\tau .}

\subsection*{Summary}

All subsequent structures --- triadic confirmation, the binary alphabet,
the \ensuremath{1\to 4} multiplier, FCC geometry, the conjugate field, the
[[4,2,2]] structural correspondence, the identification of the
Planck length and Planck time, the derivation of \ensuremath{c_{\mathrm{TMS}}} = \ensuremath{1/\sqrt{2},} and the
wave equation --- are derived from these axioms and the Principle of
Generative Economy without additional assumption.

\subsection*{1.1 Maximum Entropy Stochastic Field \ensuremath{(U_{\mathrm{sh}})}}

The ground state of the TMS is \ensuremath{U_{\mathrm{sh}}} as defined in Axiom 0 --- an
irreducible, algorithmically incompressible binary flux. It is the
primordial field of potential from which all structured information is
actualized. \ensuremath{U_{\mathrm{sh}}} acts on the informational layer --- the latent bits
populating the interstices of \ensuremath{M\alpha .} It does not affect the agents, which
remain immutable substrate.

\subsection*{1.2 Agent Manifold \ensuremath{(M\alpha )}}

A 2D hexagonal array of spherical primitives \ensuremath{(\alpha ).} Agents are the
immutable geometric substrate of the TMS. They do not change, move, or
degrade. By Axiom 2 and the Principle of Generative Economy, the sphere
is the unique primitive --- zero orientation parameters, minimal
surface-to-volume ratio. Any other primitive introduces orientation
degrees of freedom inconsistent with the Principle of Generative
Economy.

\subsection*{1.3 Touch Vectors}

Each agent \ensuremath{\alpha_{i}} maintains six touch vectors connecting it to its six
neighbors. Touch vectors are the channels of informational communication
between agents. A channel is shared infrastructure --- the same vector
can carry signals relating to different bits --- but each signal is
unique to the specific bit it describes. Touch vectors point either
laterally within a layer or forward into the next layer. There is no
backward propagation. Each signal carried by a touch vector inherits the
polarization \ensuremath{\chi} \ensuremath{\in} \{\ensuremath{-1,} +1\} of the confirmation event that produced it.

\subsection*{1.4 Bits, Confirmation, and the Binary Alphabet}

Bits (B) exist in the interstices of \ensuremath{M\alpha} in a latent state, drawn from
\ensuremath{U_{\mathrm{sh}}.} By the Principle of Generative Economy, a bit is exactly the size
of the triangular interstice formed by three touching spheres --- the
unique scale at which a bit can be unambiguously localized with complete
angular closure. Actualization --- the transition of latent B to
confirmed \ensuremath{\bar{B}} --- requires triadic agreement among agents.

A confirmed bit carries a binary value:

\emph{\ensuremath{\bar{B}} \ensuremath{\in} \{0, 1\}}

\ensuremath{\bar{B}} = 1 denotes a confirmed presence --- a positive informational result.
\ensuremath{\bar{B}} = 0 denotes a confirmed absence --- a negative informational result.
Both are stabilized outcomes pulled from \ensuremath{U_{\mathrm{sh}}} by triadic confirmation;
both lock the interstice into the causal record; both render a confirmed
tetrahedral apex at the next layer (Section II). The distinction between
the two is not whether the interstice is locked (every confirmed
interstice is locked) but what computational content is locked there.

The latent state B in \ensuremath{U_{\mathrm{sh}}} --- neither 0 nor 1, but the algorithmically
incompressible superposition of both --- is ontologically distinct from
the confirmed zero \ensuremath{\bar{B}} = 0. A latent bit is unmeasured; a confirmed zero
is the measured outcome of an absence. This trinary distinction ---
latent, confirmed-zero, confirmed-one --- is the substrate's irreducible
computational state space.

The signed representation:

\emph{\ensuremath{\chi} \ensuremath{\equiv} \ensuremath{2\bar{B}} \ensuremath{-} 1 \ensuremath{\in} \{\ensuremath{-1,} +1\}}

is referred to as the polarization of the confirmed bit and is used in
subsequent derivations.

Only bits actualize and deactualize; agents do not. The mechanism by
which incoming signals bias the resolution of a confirmation toward \ensuremath{\bar{B}} =
1 or \ensuremath{\bar{B}} = 0 is developed in Volume 1, Paper 2. This paper treats the
binary alphabet as an ontological primitive established by the
interaction of triadic confirmation with the maximum-entropy substrate.

\subsection*{1.5 The Flop}

A flop is the atomic unit of state change in the TMS. It is the complete
event sequence initiated by a triadic confirmation: the locking of a bit
\ensuremath{\bar{B}} with binary value, the simultaneous generation of 12 outgoing touch
vector signals carrying the polarization of the confirmation, the
consumption of those signals as 4 synthetic bits, and the resulting
Z-displacement to a new layer. A flop is indivisible --- all sub-events
occur simultaneously as a single act. There is no sub-flop time. The
discrete time parameter \ensuremath{\tau} counts flops. All rates in the TMS ---
confirmation rate, synthesis rate, propagation speed --- are expressions
of this single underlying tempo.

These primitives and their interactions define the complete ontology of
the TMS. All subsequent structures are derived from them without
additional assumption.

\section*{II. FCC Necessity}
\addcontentsline{toc}{section}{II. FCC Necessity}

\subsection*{2.1 The Triadic Minimum}

\begin{theorem}[Triadic Minimum]
3 agents is the necessary and sufficient
minimum for bit actualization in \ensuremath{M\alpha .}
\end{theorem}

\begin{proof}
The confirmation substrate must constitute a
closed simplicial cover of the plane --- every interstice fully bounded
with complete angular closure. No open edges are permitted; a bit cannot
be localized without a closed boundary.

Case 1: 2 agents is insufficient. Two agents define a 1-simplex. A
1-simplex has no interior and subtends less than \ensuremath{360^{\circ}} regardless of
arrangement. The enclosure is always open. No bit can be localized.
2-agent confirmation is geometrically impossible.

Case 2: 4 agents is non-minimal. Four agents per interstice produces a
closed simplicial cover --- the Cartesian square tiling, \ensuremath{\bar{B}} = \ensuremath{4\bar{\alpha}_{2}.} This
is geometrically valid but violates the Principle of Generative Economy:
it achieves complete angular closure at one degree of complexity greater
than necessary.

Case 3: 3 agents is necessary and sufficient. Three agents define a
2-simplex. A 2-simplex achieves complete angular closure --- \ensuremath{360^{\circ}} ---
with no redundancy. Removing any one agent opens the boundary. Adding
any agent violates the Principle of Generative Economy. The triangular
confirmation geometry tiles the plane completely as a closed simplicial
cover, producing the hexagonal lattice as consequence.

Therefore \ensuremath{\bar{B}} = \ensuremath{3\bar{\alpha}_{2}} is the unique solution selected by the Principle of
Generative Economy. The hexagonal lattice is not assumed --- it follows
from the triadic minimum.

\end{proof}

\begin{equation*}
\bar{B} = 3\bar{\alpha}_{2} (2D)
\end{equation*}

\subsection*{2.2 Hexagonal Packing}

Given the sphere as the unique geometric primitive under Axiom 2, the
densest packing of equal discs in 2D is the hexagonal arrangement (Thue,
1890; rigorously proved Fejes Tóth, 1940). This is the unique
maximum-density solution and defines the structure of \ensuremath{M\alpha .} It is the
direct consequence of the triadic minimum --- the hexagonal lattice is
what a closed simplicial cover of circles necessarily produces.

\subsection*{2.3 Causal Anchoring and Z-Displacement}

When a bit is confirmed in layer \ensuremath{S_{0},} the confirming agents and their
relationships constitute the causal record of that event. By Axiom 4,
this record cannot be overwritten. The \ensuremath{1\to 4} multiplier (Section IV)
generates synthetic bits that must be resolved at new positions. Three
candidate directions exist:

Lateral expansion within \ensuremath{M\alpha} is geometrically impossible. By Axiom 2 and
the Principle of Generative Economy, \ensuremath{M\alpha} is hexagonally close-packed ---
the unique maximum-density arrangement of equal spheres in 2D (Thue,
1890). No unoccupied interstice exists at the same layer into which a
new agent-scale sphere can be inserted without overlapping an existing
agent. Lateral expansion violates Axiom 2.

Oblique displacement introduces a free parameter. Any displacement at an
angle other than perpendicular to \ensuremath{M\alpha} requires specifying that angle ---
an orientation parameter with no geometric basis in the existing
structure. This violates the Principle of Generative Economy.

Orthogonal displacement is geometrically forced. When a new sphere of
identical radius settles into the triangular interstice formed by three
mutually touching spheres of \ensuremath{M\alpha ,} its center comes to rest at a uniquely
determined height above the plane --- fixed entirely by the contact
geometry of equal spheres. This height defines the perpendicular
direction. The third dimension is not imported from a pre-existing
Euclidean container. It is the unique direction produced by
sphere-on-sphere contact in the triangular interstice. Orthogonality is
derived, not assumed.

The displacement therefore defines the Axis of Causal Succession:

\begin{equation*}
S_{0} \to S_{1} : \Delta z = 1
\end{equation*}

Each new layer \ensuremath{S_{n}} is the physical record of a confirmed state. This
directional necessity is the geometric origin of flop time \ensuremath{\tau .}

\subsection*{2.4 FCC Selection over HCP}

When a second layer \ensuremath{S_{1}} forms above \ensuremath{S_{0},} it occupies the interstices of
the first --- position set B over position set A. The two candidate
stacking sequences for a third layer are ABABAB (hexagonal close-packed,
HCP) and ABCABC (face-centered cubic, FCC).

To evaluate these candidates, we must first define what constitutes a
distinct position in the causal stack.

\begin{definition}[Causal Position]
A position in the TMS causal stack is
defined by its complete touch vector connectivity --- the full set of
geometric relationships between that node and every node in all prior
layers. Two positions are causally identical if and only if their
connectivity patterns are geometrically indistinguishable to any
downstream signal.
\end{definition}

This definition is forced by Axiom 4. The causal record is a record of
signal history --- what arrived, from where, and in what geometric
relationship. If two positions share identical connectivity patterns,
then a signal reaching either position carries no geometric information
about which layer it originated from. The causal record collapses.

\begin{theorem}[FCC Necessity]
Given Definition (Causal Position) and Axiom 4,
the unique permitted stacking sequence is ABCABC.
\end{theorem}

\begin{proof}
In the ABABAB sequence, \ensuremath{S_{2}} returns to A positions. The touch
vector connectivity pattern of an A-position node in \ensuremath{S_{2}} is geometrically
congruent to that of the corresponding A-position node in \ensuremath{S_{0}} --- both
sit directly above B-position nodes in an intervening layer, with
identical angular relationships to their neighbors. By Definition
(Causal Position), these nodes are causally identical. A signal
propagating through \ensuremath{S_{1}} into \ensuremath{S_{2}} cannot geometrically distinguish whether
it has arrived at a new causal record or returned to \ensuremath{S_{0}.} The causal
stack becomes periodic with period 2 --- a closed causal loop. This
violates Axiom 4: the confirmed state of \ensuremath{S_{0}} is effectively overwritten
by the indistinguishability.

In the ABCABC sequence, \ensuremath{S_{2}} occupies C positions --- geometrically
distinct from both A and B. The touch vector connectivity of every
C-position node is unique: it is above a B-position node in \ensuremath{S_{1},} which is
itself above an A-position node in \ensuremath{S_{0},} producing an asymmetric
connectivity pattern that distinguishes \ensuremath{S_{2}} from every prior layer. No
two layers share a connectivity pattern. The causal record is
monotonically distinguishable at every layer.

HCP is excluded by causal anchoring. FCC is the unique stacking sequence
consistent with Axiom 4.

\end{proof}

\begin{corollary}[Sublattice Bifurcation at Recursive Levels]
The FCC structure derived here applies at the substrate level. At
each hierarchical level k \ensuremath{\geq} 1 produced by recursive coarse-graining
(Paper 3 \S\,4.5), the meta-FCC sublattice further splits into two
interpenetrating FCC sublattices related by spatial inversion through
the (1,1,1)-direction offset. The full bifurcated structure that
results, and the empirical verification of this bifurcation through
\ensuremath{L_{\mathrm{box}}} = 12, is derived in Paper 5 \S\,6.1.3 and Appendix E. The
substrate-level FCC of this paper is the level-0 case of the
sublattice-bifurcation recursion that operates at every level k \ensuremath{\geq} 1.
\end{corollary}

\subsection*{2.5 Triangular to Tetrahedral Interstices}

The ABCABC stacking sequence converts triangular interstices bounded by
3 agents into tetrahedral interstices bounded by 4 agents. The
confirmation geometry is preserved across the dimensional transition:

\begin{equation*}
\bar{B} = 3\bar{\alpha}_{2} (2D)
\end{equation*}

\begin{equation*}
\bar{B} = 4\bar{\alpha}_{3} (3D)
\end{equation*}

The triadic rule does not change. The geometry scales naturally from
triangle to tetrahedron without introducing any new mechanism. The
binary alphabet \ensuremath{\bar{B}} \ensuremath{\in} \{0, 1\} of Section 1.4 is preserved unchanged: both
2D and 3D confirmation events resolve to a binary value.

The tetrahedral confirmation unit --- four agents enclosing one bit with
one of two possible values --- emerges as the fundamental 3D structure.
Its informational dynamics are derived in Section IV.

\section*{III. The Fundamental Length Scale}
\addcontentsline{toc}{section}{III. The Fundamental Length Scale}

With the FCC geometry established, the system possesses a unique length
scale --- the agent diameter. This section shows that scale is not a
free parameter but is forced to equal the Planck length, and derives the
Planck time as a necessary consequence.

\subsection*{3.1 Definition}

\begin{definition}[Fundamental length]
Let \ensuremath{\ell_{0}} denote the diameter of the
spherical primitive \ensuremath{\alpha .} This is the unique length scale appearing in the
definition of \ensuremath{M\alpha .}
\end{definition}

\subsection*{3.2 Planck Length Identification}

\begin{theorem}[Planck Length Identification]
The fundamental agent spacing \ensuremath{\ell_{0}}
is identical to the Planck length \ensuremath{\ell_{p}.} This identification carries no
free parameters.
\end{theorem}

\begin{proof}
The 2D manifold \ensuremath{M\alpha} is the complete pre-geometric substrate --- no
structure exists below it. Any measurement of distance in the emergent
3D geometry must ultimately resolve to positions of agents or
interstices within \ensuremath{M\alpha .} Two distinct agents cannot occupy positions
closer than \ensuremath{\ell_{0},} as their spheres would overlap, violating Axiom 2.
Therefore \ensuremath{\ell_{0}} is the minimum resolvable length in the emergent 3D
geometry. Below \ensuremath{\ell_{0},} no geometric distinctions exist --- only the 2D
substrate.

The Planck length \ensuremath{\ell_{p}} = \ensuremath{\sqrt{\hbar G/c^{3}}} is the scale below which classical
spacetime geometry ceases to be meaningful. The TMS provides the
geometric reason: below \ensuremath{\ell_{0},} the emergent 3D geometry does not exist. The
identification

\begin{equation*}
\ell_{0} \equiv \ell_{p}
\end{equation*}

is therefore forced by the ontology. It is not an external input --- it
is the statement that the pre-geometric substrate sets the boundary of
geometry. No free parameter remains.

\end{proof}

\begin{remark}[On the Status of \ensuremath{\hbar} and G]
The identification \ensuremath{\ell_{0}} \ensuremath{\equiv} \ensuremath{\ell_{p}} fixes the substrate's minimum length,
and \S\,3.3 fixes its tick \ensuremath{\tau_{\mathrm{flop}}}; from these the speed \ensuremath{c} is
derived. It is worth stating precisely what this leaves, because the status of
\ensuremath{\hbar} and G is a point of principle, not a deferred calculation.

The substrate is combinatorial: it produces dimensionless quantities --- counts,
ratios, angles (the coordination number 12, the subshell capacities 2, 6, 10, 14,
the noble-gas atomic numbers, the fermionic sign, the reconstructed angles of
Paper 5 \S\,4.1). A dimensionful constant cannot be produced by dimensionless
structure alone; it requires a unit, which is scale-setting data, not a number the
combinatorics can output. The three constants \ensuremath{\{c, \hbar, G\}} form an invertible
basis for the three mechanical dimensions (mass, length, time) --- fixing them is
exactly the choice of Planck units --- so any theory expressed in dimensionful
units must be handed a scale for each dimension it cannot supply intrinsically.

The TMS accounting is favorable and exact. The substrate supplies \emph{two} of
the three scales from its own structure: the length \ensuremath{\ell_{p}} (lattice geometry)
and the time \ensuremath{\tau_{\mathrm{flop}}} (the update tick). Their ratio gives the third
mechanical quantity of the length--time sector, the speed \ensuremath{c} (\S\,3.3), so \ensuremath{c}
is derived rather than imported. What remains is the mass/action sector. The
substrate's fundamental quantum is a \emph{bit} --- dimensionless information ---
not an action; a combinatorial substrate therefore has no intrinsic action scale,
and exactly one dimensionful input must be supplied to reach physical units. That
single input is what \ensuremath{\hbar} carries. Newton's constant is then not independent:
given \ensuremath{\ell_{p}}, \ensuremath{c}, and \ensuremath{\hbar}, the Planck relation
\ensuremath{\ell_{p} = \sqrt{\hbar G/c^{3}}} inverts to \ensuremath{G = c^{3}\ell_{p}^{2}/\hbar}, so G is
\emph{derived}. The framework thus takes exactly one dimensionful scale-setting
input, \ensuremath{\hbar}, and yields \ensuremath{c} and G --- one import, not two open constants, and
the minimum any theory stated in dimensionful units can require.

This sharpens, rather than weakens, the no-free-parameters claim. TMS has no free
\emph{dimensionless} parameters: every dimensionless quantity is fixed by the
axioms. It has exactly one dimensionful input, \ensuremath{\hbar}, which sets the scale at
which the dimensionless structure is read --- and even that does not float freely,
since it fixes \ensuremath{c} and G with it. Asking the framework to output the numerical
\emph{value} of \ensuremath{\hbar} in SI would be a dimensional category error, equivalent to
asking it to output the metre or the second: that value is the definition of the
unit system, not a fact the combinatorics decides. This is the same footing on
which the Standard Model fixes its group-theoretic structure from first principles
while taking its dimensionful couplings from experiment.
\end{remark}

\subsection*{3.3 Flop Time}

\begin{corollary}[Flop Time]
The substrate has two intrinsic dimensionful scales --- the minimum length
\ensuremath{\ell_{p}} (\S\,3.2) and the flop time \ensuremath{\tau_{\mathrm{flop}}}, the duration of a single
update. The propagation speed is the ratio of the two, and the flop time is
identified with the Planck time:
\end{corollary}

\begin{equation*}
c = \frac{1}{\sqrt{2}}\,\frac{\ell_{p}}{\tau_{\mathrm{flop}}}, \qquad
\tau_{\mathrm{flop}} = t_{p}
\end{equation*}

\begin{proof}
The flop is the substrate's irreducible tick: there is no sub-flop time
(\S\,II), and the discrete time parameter \ensuremath{\tau} counts flops. The flop time
\ensuremath{\tau_{\mathrm{flop}}} is therefore an intrinsic primitive of the dynamics, not a
quantity constructed from other scales. Together with the intrinsic length
\ensuremath{\ell_{p}}, it fixes a speed: a signal advances one face-diagonal step per flop,
giving \ensuremath{c = (1/\sqrt{2})\,\ell_{p}/\tau_{\mathrm{flop}}} in lattice units (the
\ensuremath{1/\sqrt{2}} is the FCC face-diagonal factor derived in \S\,VI). The speed of
light is thus an \emph{output} of the substrate --- the lattice's own
steps-per-tick rate --- not an external input. The wave equation of \S\,VI
confirms the identification: the physical wave speed from the discrete dynamics
equals \ensuremath{c} precisely when the flop time is the Planck time, \ensuremath{\tau_{\mathrm{flop}}}
= \ensuremath{t_{p}} = \ensuremath{\sqrt{\hbar G/c^{5}}}. One flop equals exactly one Planck time.
\end{proof}

\subsection*{3.4 Agent Density}

\begin{corollary}[Agent Density]
The agent density in \ensuremath{M\alpha} is \ensuremath{2/(\sqrt{3}} \ensuremath{\cdot} \ensuremath{\ell^{2}_{p})} per
unit area. The FCC lattice inherits the same fundamental spacing \ensuremath{\ell_{p}}
between agent centers.
\end{corollary}

\section*{IV. The Triadic Actualization Rule and the \ensuremath{1\to 4} Multiplier}
\addcontentsline{toc}{section}{IV. The Triadic Actualization Rule and the \ensuremath{1\to 4} Multiplier}

\subsection*{4.1 The Triadic Constraint}

The minimum boundary condition for bit actualization in \ensuremath{M\alpha} is a
2-simplex (triangle) formed by three agents, as derived in Section II.
Each agent has six touch vectors. Each bit costs exactly three unique
signals --- one from each vertex of the confirming triangle. An agent
may participate in up to six confirmations (one per neighbor pair), but
each confirmation consumes three unique signals. No signal describes
more than one bit.

\begin{equation*}
\bar{B} = 3\bar{\alpha}_{2}
\end{equation*}

\subsection*{4.2 The Multiplier Proof}

Given a primary confirmation event \ensuremath{S_{0}} involving agents \{\ensuremath{\bar{\alpha}_{1},} \ensuremath{\bar{\alpha}_{2},} \ensuremath{\bar{\alpha}_{3}}\}:

\begin{itemize}
\item
  Each confirming agent has 6 touch vectors
\item
  2 vectors per agent are locked into the confirmation triangle
\item
  4 vectors per agent carry the signal outward
\item
  Total outgoing signals: 3 \ensuremath{\times} 4 = 12
\end{itemize}

Since each bit costs exactly three unique signals:

\begin{equation*}
12 \div 3 = 4 \text{synthetic} \text{bits}
\end{equation*}

This division is exact by necessity. The hexagonal geometry is the
unique 2D packing where 6 neighbors per agent produces a clean integer
result under the triadic constraint. No remainder is geometrically
permitted --- any configuration that would produce one is not a valid
confirmation geometry.

\subsection*{4.3 Informational Closure and Polarization Transport}

Every confirmation event is informationally complete within a single
flop. All 12 outgoing signals are fully consumed by the 4 synthetic
bits. Nothing is stored, deferred, or lost. The TMS is informationally
closed by construction --- a direct consequence of Axiom 3. Closure here
is of signal count and content: every signal is consumed, none created or
destroyed (Axiom 3). What the signed aggregate does not preserve is
resolvable distinction --- three incoming polarizations collapse to one
signed outcome --- and this per-flop loss of distinction, not of content,
is the geometric origin of the irreversibility developed in Paper 5.

The 12 outgoing signals carry the polarization of their source
confirmation. A \ensuremath{\bar{B}} = 1 confirmation emits 12 signals carrying \ensuremath{\chi} = +1; a \ensuremath{\bar{B}}
= 0 confirmation emits 12 signals carrying \ensuremath{\chi} = \ensuremath{-1.} The signal count is
identical in both cases --- informational closure operates on count
alone --- but the polarization content differs.

The \ensuremath{1\to 4} multiplier is content-preserving: each downstream synthetic bit
position accumulates the polarizations of its 3 incoming signals as a
signed aggregate. The resulting polarization landscape across the
synthetic bit positions of \ensuremath{S_{n+1}} biases the subsequent resolution of
those positions toward \ensuremath{\bar{B}} = 1 or \ensuremath{\bar{B}} = 0 against the unbiased stochastic
kick from \ensuremath{U_{\mathrm{sh}}.} The propagation dynamics of this polarization landscape,
the per-flop biasing mechanism, and the structural correspondence to the
Born rule are developed in Volume 1, Paper 2.

The informational closure of each confirmation event, combined with the
\ensuremath{1\to 4} multiplier, generates a structural surplus that cannot be resolved
within the existing 2D layer --- necessitating the Z-displacement and
FCC geometry derived in Section II. The stability properties of the
resulting tetrahedral structure are examined in Section V.

\section*{V. The [[4,2,2]] Structural Correspondence and Fault-Tolerant Propagation}
\addcontentsline{toc}{section}{V. The [[4,2,2]] Structural Correspondence and Fault-Tolerant Propagation}

\subsection*{5.1 The Stabilizer Mapping}

The [[4,2,2]] quantum error-detection code encodes 2 logical
qubits in 4 physical qubits with code distance 2. Its stabilizer group
is generated by:

\begin{equation*}
S_{x} = X_{1}X_{2}X_{3}X_{4}
\end{equation*}

\begin{equation*}
S_{Z} = Z_{1}Z_{2}Z_{3}Z_{4}
\end{equation*}

A state lies in the code space when both stabilizers evaluate +1. Any
single-qubit error anticommutes with at least one stabilizer, producing
a measurable syndrome (the stabilizer evaluates to \ensuremath{-1).} The
two-dimensional code space encodes 2 independent logical qubits --- 4
distinguishable logical states.

The tetrahedral confirmation unit of the TMS maps onto this structure as
follows:

\begin{itemize}
\item
  4 physical qubits \ensuremath{\to} 4 agents \{\ensuremath{\bar{\alpha}_{1},} \ensuremath{\bar{\alpha}_{2},} \ensuremath{\bar{\alpha}_{3},} \ensuremath{\bar{\alpha}_{4}}\} of the tetrahedral
  interstice.
\item
  Logical qubit 1 \ensuremath{\to} the value of the confirmed bit in the current plane:
  \ensuremath{\bar{B}} \ensuremath{\in} \{0, 1\}.
\item
  Logical qubit 2 \ensuremath{\to} the inheritance of that value through the causal
  stack --- the integrity of the bit's history preserved across prior
  layers.
\end{itemize}

Every confirmed bit therefore carries two independent informational
degrees of freedom: its current value and its causal history. The 4
logical states of the [[4,2,2]] code space correspond exactly to
the 4 possible combinations of these two qubits. The encoding is
faithful.

Single-agent disruptions caused by \ensuremath{U_{\mathrm{sh}}} fall into two algebraically
distinct types, each detected by a different stabilizer:

\begin{itemize}
\item
  Bit-flip-type disruption (X-type error). A disturbance that would
  alter the value contribution of a single agent. This anticommutes with
  the Z-stabilizer, producing \ensuremath{S_{Z}} = \ensuremath{-1.}
\item
  Phase-type disruption (Z-type error). A disturbance that would corrupt
  the inheritance contribution of a single agent. This anticommutes with
  the X-stabilizer, producing \ensuremath{S_{x}} = \ensuremath{-1.}
\end{itemize}

Both syndrome flavors are checked at every flop. The confirmation
mechanism that pulls a bit from \ensuremath{U_{\mathrm{sh}}} is the triadic actualization rule of
Section II; the [[4,2,2]] stabilizers do not perform
confirmation but rather verify the integrity of already-confirmed bits
against subsequent disruption. The dynamic-code framing of recent work
on stabilizer codes emerging from spacetime circuits (Aitchison \& Béri,
2025) provides a parallel formal context for the structural
correspondence developed here.

\subsection*{5.2 The Restoring Force}

When \ensuremath{U_{\mathrm{sh}}} disrupts a confirmed bit \ensuremath{\bar{B}} --- flipping its value or corrupting
its causal history --- the 4 agents of the tetrahedral unit are
unchanged. They remain in place as immutable substrate. The confirmation
relationship between them is violated, producing a detectable parity
syndrome.

The synthetic bits produced at the current flop --- the same flop at
which the violation is detected --- repair the missing or corrupted
confirmation, restoring parity. Detection and repair are a single atomic
act, a direct consequence of informational closure. The system cannot
defer repair any more than it can defer signal consumption. Everything
resolves within the flop.

Both logical qubits are detected and repaired within the same flop. A
bit-flip-type violation \ensuremath{(S_{Z}} = \ensuremath{-1)} and a phase-type violation \ensuremath{(S_{x}} = \ensuremath{-1)}
are independent failure modes, protected independently. In every flop,
the TMS simultaneously maintains the present state, the binary value of
every confirmed bit, and the integrity of the entire causal record. The
system is not merely propagating forward --- it is actively preserving
the full informational content of its history as a single geometric
motion.

\subsection*{5.3 Code Distance and Fault-Tolerant Propagation}

The [[4,2,2]] code has distance 2 --- any single error is
detectable. In TMS terms:

\begin{itemize}
\item
  A single bit-flip-type disruption on one agent violates \ensuremath{S_{Z}} while
  leaving \ensuremath{S_{x}} intact --- detectable.
\item
  A single phase-type disruption on one agent violates \ensuremath{S_{x}} while leaving
  \ensuremath{S_{Z}} intact --- detectable.
\item
  The tetrahedral geometry guarantees that no single-agent disruption
  can leave both stabilizers intact --- code distance 2 follows from the
  geometry alone.
\end{itemize}

The full stabilizer algebra --- both stabilizers, both logical qubits,
the assignment of value and history to the two logical degrees of
freedom, and the code distance --- is derived from TMS primitives. The
mapping constitutes a structural correspondence. The framework is also
consistent with the broader holographic quantum error correction
tradition (Pastawski, Yoshida, Harlow, \& Preskill, 2015; El Morsalani,
2025) that has informed the use of stabilizer codes in
emergent-spacetime contexts. \hfill\ensuremath{\square}

The [[4,2,2]] structural correspondence guarantees that the TMS
substrate is fault-tolerant against stochastic interference from \ensuremath{U_{\mathrm{sh}}.}
Confirmed bits propagate across the lattice with their values and
histories intact, because the error-detection property of each
tetrahedral unit preserves both at every hop. This establishes the TMS
lattice as a stable propagation medium for both the confirmation field
and the polarization landscape carrying binary content.

The discrete dynamics of this fault-tolerant lattice, when taken to the
continuum limit, reproduce the standard wave equation --- derived in
Section VI.

The structural correspondence developed here applies at the substrate
level. The same [[4,2,2]] stabilizer structure reappears at
every hierarchical level by axiom scale-independence (Paper 3 \S\,4.7),
with one structural addition at level k \ensuremath{\geq} 1: the meta-FCC splits into
two interpenetrating sublattices A and B (Paper 3 \S\,4.5; Paper 5 \S\,6.1.3),
each running independent [[4,2,2]] dynamics. The level-1
verification is empirical (Paper 5 Appendix E, Stages 5 and 6). The
cross-level commutator structure between the A-sublattice and
B-sublattice stabilizers at level k provides the substrate-level origin
of non-Abelian gauge closure at level k + 1 (Paper 5 \S\,8.2).

\subsection*{5.4 Scale-Invariant Stabilizer Null-Space}

The [[4,2,2]] stabilizer structure of \S\,5.1--\S\,5.3, applied to an
FCC patch of finite extent \ensuremath{L_{\mathrm{box}}} at any hierarchical level, has a
precise count of stabilizer-satisfying configurations: the dimension of
the null space over \ensuremath{F_{2}} of the tetrahedron-to-site parity matrix. We
state this as a theorem.

\begin{theorem}[Scale-Invariant Stabilizer Null-Space]
For an
FCC patch of edge \ensuremath{L_{\mathrm{box}}} at any hierarchical level, with
[[4,2,2]] stabilizers imposed at every tetrahedron in the patch,
the dimension of the stabilizer-satisfying configuration space over \ensuremath{F_{2}}
is k = 3n, where n = \ensuremath{(L_{\mathrm{box}}} \ensuremath{-} 1)/2 at level 1 and analogously at higher
levels (the level-local lattice depth). The number of distinct
stabilizer-satisfying configurations is therefore \ensuremath{2^{(3n)},} and this
scaling law is a structural property of the FCC + [[4,2,2]]
dynamics that is independent of hierarchical level.
\end{theorem}

\begin{proof}
The [[4,2,2]] stabilizers at each tetrahedron act
linearly on the \ensuremath{F_{2}-\text{valued}} polarization state of the four agents in the
tetrahedron. Their joint action on a patch of N sites with T tetrahedra
is a linear map \ensuremath{F_{2}^{N}} \ensuremath{\to} \ensuremath{F_{2}^{T}} whose kernel dimension is the
satisfying configuration count. The geometric argument that this kernel
dimension equals 3n at every level is given in schematic form by the FCC
cubic-volume scaling: the unit cubic cell of side 2 in level-local
lattice units contributes 3 degrees of freedom to the null space (one
per axis-pair of the cubic cell), and a patch of cubic-cell depth n
contains \ensuremath{n^{3}} unit cells along with their bounding stabilizer relations,
giving 3n dimensions of freedom from the cubic-cell-edge degrees and \ensuremath{(n^{3}}
\ensuremath{-} 3n) parity constraints from the tetrahedron-interior stabilizers. The
result is k = 3n. The empirical verification of this law at substrate
level (n = 1 through \ensuremath{L_{\mathrm{box}}} extents matching the patch sizes) and at
levels 1 and 2 (Paper 5 Appendix E, Stages 5 and 6) confirms k = 3n
through n = 6 across all levels tested.

\end{proof}

\begin{remark}
The scale-invariance of the k = 3n law is the
structural reason the recursion-preservation theorem of Paper 5 \S\,2.5 has
its full force: the [[4,2,2]] dynamics at every level have the
same combinatorial structure when measured in level-local units, because
the null-space scaling depends only on the FCC patch geometry, not on
the absolute scale.
\end{remark}

\section*{VI. The Wave Equation}
\addcontentsline{toc}{section}{VI. The Wave Equation}

\subsection*{6.1 The Confirmation Field}

Let \ensuremath{\varphi (\bar{r},} \ensuremath{\tau )} denote the local confirmation density --- the scalar field
representing the density of actualized bits \ensuremath{\bar{B}} at any point in the
lattice:

\begin{equation*}
\varphi (\bar{r}, \tau ) = \text{confirmed} \text{bits} in \text{region} / \text{total} \text{available} \text{interstices}
in \text{region}
\end{equation*}

This field counts confirmation events without regard to their binary
value. It is purely informational and purely geometric. It requires no
additional physical assumptions.

\subsection*{6.2 The Discrete Laplacian on the FCC Lattice}

Each node in the FCC lattice has 12 nearest neighbors. The discrete
Laplacian is defined as:

\begin{equation*}
\Delta_{\mathrm{FCC}} \varphi_{i} = \Sigma (\varphi_{j} - \varphi_{i}), j \in \text{neigh}(i)
\end{equation*}

This operator measures the difference between the confirmation density
at node i and its neighbors --- the local informational gradient across
the lattice.

\subsection*{6.3 The Conjugate Field and Second-Order Dynamics}

\begin{lemma}[Field Conjugacy]
The confirmation field \ensuremath{\varphi} and the
forward-propagating excitation field \ensuremath{\Pi} constitute a co-generated
conjugate pair, arising necessarily from the interaction between the \ensuremath{1\to 4}
cascade and \ensuremath{U_{\mathrm{sh}}} regulation.
\end{lemma}

\begin{proof}


Step 1 --- The 2D case produces first-order closure. In the 2D manifold
\ensuremath{M\alpha} without Z-displacement, all signals generated by a confirmation event
are consumed within the same layer. The state at any layer is fully
described by \ensuremath{\varphi} alone. No forward-propagating field persists. The update
rule is first-order by construction --- confirmation density at the next
step depends only on confirmation density at the current step.

Step 2 --- The \ensuremath{1\to 4} cascade generates exponential confirmation pressure.
In the 3D FCC lattice, each confirmation event at layer \ensuremath{S_{n}} generates 12
outgoing signals, consumed as 4 synthetic confirmations at \ensuremath{S_{n+1}.} Each of
those 4 confirmations generates 12 signals, consumed as 16 confirmations
at \ensuremath{S_{n+2}.} Left unchecked, this produces an exponential cascade:

\emph{1 \ensuremath{\to} 4 \ensuremath{\to} 16 \ensuremath{\to} 64 \ensuremath{\to} \ldots{}}

This cascade is not a pathology --- it is the direct geometric
consequence of informational closure applied layer-to-layer under FCC
connectivity. Every signal is consumed. Every consumption generates new
signals. The system is structurally amplifying.

Step 3 --- \ensuremath{U_{\mathrm{sh}}} regulation produces dynamic stability. \ensuremath{U_{\mathrm{sh}}} acts on the
latent bit population at every layer, disrupting confirmations
stochastically and returning bits to the unactualized field. This
regulation is the only mechanism that competes with the \ensuremath{1\to 4} cascade.

The cascade generically outpaces disruption because amplification is
deterministic and geometric while disruption is random. At each layer,
the \ensuremath{1\to 4} mechanism advances with fixed multiplicity; \ensuremath{U_{\mathrm{sh}}} disrupts with
probability distributed across the full bit population. On average, the
cascade stays ahead --- not because the system reaches equilibrium, but
because geometric amplification is structurally faster than
uncoordinated stochastic interference.

The [[4,2,2]] fault tolerance (Section V) reinforces this at the
local level. Each tetrahedral confirmation unit detects and repairs
single disruptions within the same flop, preserving causal coherence
node by node. As the structure grows, its surface area exposed to \ensuremath{U_{\mathrm{sh}}}
increases --- but so does the interior volume protected by the causal
stack. The propagating wave is the coherent interior structure advancing
faster than disruption can reach it.

This is dynamic stability, not static equilibrium. The system does not
settle --- it persists by continuing to move forward, in the same way a
bicycle is stable under motion but not at rest. A sufficiently large
simultaneous disruption --- a true stochastic black swan event
collapsing enough bits at once to exceed the current flop's resolution
capacity --- would constitute a system-wide failure. Such a failure
produces no downstream causal layer and therefore leaves no record
within the stack. By Axiom 4, only what is confirmed is recorded. A
collapsed system is unobservable from within by construction.

The propagating residual \ensuremath{\Pi} is therefore the net forward-propagating
confirmation pressure in a dynamically stable instance --- the touch
vector field carrying well-defined inherited excitation forward across
layers under the constraint of Axiom 3.

Step 4 --- Signal conservation forces the coupling equations. At any
layer \ensuremath{S_{n},} the system state requires two quantities for complete
specification:

\begin{itemize}
\item
  \ensuremath{\varphi (S_{n})} --- the confirmation density: what actually locked at this layer
\item
  \ensuremath{\Pi (S_{n})} --- the inherited excitation pressure: what is being driven
  forward into \ensuremath{S_{n+1}}
\end{itemize}

By Axiom 3, signals are conserved across flops. The rate of change of
confirmation density is exactly the inherited excitation pressure --- no
signal is created or destroyed in the transfer:

\begin{equation*}
\partial \varphi / \partial \tau = \Pi
\end{equation*}

The rate of change of excitation pressure is governed by the local
confirmation gradient across the FCC lattice:

\begin{equation*}
\partial \Pi / \partial \tau = c^{2} \Delta_{\mathrm{FCC}} \varphi
\end{equation*}

These equations are not assumed. They are the unique expressions of
Axiom 3 applied to a two-quantity state under FCC connectivity.

\end{proof}

\begin{remark}[Uniqueness of the Coupling]
The second coupling
equation \ensuremath{\partial \Pi /\partial \tau} = \ensuremath{c^{2}\Delta_{\mathrm{FCC}}} \ensuremath{\varphi} admits no additional local term in \ensuremath{\varphi} under the
FCC touch-vector geometry. A candidate alternative of the form \ensuremath{\partial \Pi /\partial \tau} =
\ensuremath{c^{2}\Delta_{\mathrm{FCC}}} \ensuremath{\varphi} + \ensuremath{\alpha \varphi} would require that the forward-propagating excitation
pressure at a node be modified by a quantity proportional to the
confirmation density \emph{at that same node}, with no reference to its
neighbors. This is a self-amplification term: the node contributes to
its own forward pressure without communicating through the touch vector
network.
\end{remark}

Such a term is forbidden by the substrate's connectivity architecture.
By Paper 1 \S\,1.3, the touch vectors are the \emph{only} channels of
informational communication in the TMS --- every signal that updates a
quantity at a node arrives across a touch vector from a neighbor. There
is no on-site signaling channel: an agent does not communicate with
itself, and a bit cannot bias its own forward propagation independently
of its neighbors. The discrete Laplacian \ensuremath{\Delta_{\mathrm{FCC}}} is the unique
nearest-neighbor local operator consistent with this constraint --- it
sums differences \ensuremath{(\varphi_{j}} \ensuremath{-} \ensuremath{\varphi_{i})} across the 12 face-diagonal neighbors and is
identically zero for spatially uniform \ensuremath{\varphi .} An \ensuremath{\alpha \varphi} term, by contrast, is
non-zero for uniform \ensuremath{\varphi ,} which would imply that a spatially uniform
confirmation density generates forward propagation in the absence of any
spatial gradient --- propagation from no source.

Axiom 3 (signal conservation) closes this. Signals propagate; they do
not spontaneously generate. A forward pressure that responds to local
density rather than local gradients would represent a creation of
propagation content with no corresponding signal expenditure across the
touch vector network. The \ensuremath{\alpha} = 0 case is therefore the unique selection
consistent with both nearest-neighbor connectivity and signal
conservation.

The same argument forbids non-local operators (which would require
channels beyond nearest-neighbor touch vectors), fractional-derivative
operators (which would require memory of confirmation density across
multiple prior layers, violating the atomicity of the flop), and
higher-derivative spatial operators beyond \ensuremath{\Delta_{\mathrm{FCC}}} (which would require
coupling across more than nearest-neighbor distance per flop). The
discrete Laplacian \ensuremath{\Delta_{\mathrm{FCC}}} acting on \ensuremath{\varphi} alone is the unique local operator
consistent with the touch vector architecture; its coefficient \ensuremath{c^{2}} is
fixed by the propagation-speed derivation of \S\,6.5. The second coupling
equation is therefore forced in both functional form and coefficient. \hfill\ensuremath{\square}

Step 5 --- Elimination yields second-order dynamics. Eliminating \ensuremath{\Pi} from
the coupled system:

\begin{equation*}
\partial ^{2}\varphi / \partial \tau^{2} = c^{2} \Delta_{\mathrm{FCC}} \varphi
\end{equation*}

The second-order structure is not imported from classical field theory.
It is the necessary consequence of a self-amplifying confirmation
cascade regulated by a maximum-entropy stochastic field. A 2D
confirmation has no cascade, no inherited pressure, and no inertia. The
transition to 3D, forced by the \ensuremath{1\to 4} multiplier, introduces exactly the
two-quantity state that produces wave rather than diffusive behavior. \hfill\ensuremath{\square}

This result is structurally analogous to the Hamiltonian formulation of
classical field theory, where conjugate pairs (q, p) produce
second-order equations of motion via elimination of momentum (Goldstein,
Poole \& Safko, 2002). The TMS recovers this structure from the
population dynamics of a self-amplifying confirmation cascade under
stochastic regulation --- grounding the conjugate pair \ensuremath{(\varphi ,} \ensuremath{\Pi )} in the
geometric and informational architecture of the substrate itself.

\subsection*{6.4 The Discrete Wave Equation}

Each confirmation event propagates through the lattice via the touch
vector network. The rate of change of \ensuremath{\varphi_{i}} is governed by the local
gradient and the fault-tolerant restoring force established in Section
V:

\begin{equation*}
\partial ^{2}\varphi_{i} / \partial \tau^{2} = c^{2}_{\mathrm{TMS}} \Delta_{\mathrm{FCC}} \varphi_{i}
\end{equation*}

where \ensuremath{\tau} is the discrete flop time and \ensuremath{c_{\mathrm{TMS}}} is the propagation speed in
lattice spacings per flop.

\subsection*{6.5 The Propagation Speed}

\begin{theorem}[Propagation Speed]
The unique propagation speed consistent
with FCC geometry, the Planck length identification of Section III, and
the physical wave speed c is \ensuremath{c_{\mathrm{TMS}}} = \ensuremath{1/\sqrt{2}} lattice spacings per flop.
\end{theorem}

\begin{proof}
In the FCC lattice with nearest-neighbor distance \ensuremath{\ell_{p},} the 12
nearest neighbors of any node lie at vectors of the form \ensuremath{(\pm \ell_{p}/\sqrt{2},}
\ensuremath{\pm \ell_{p}/\sqrt{2},} 0) and cyclic permutations --- all directed along face
diagonals, each of magnitude \ensuremath{\ell_{p}.}

Taylor-expanding \ensuremath{\varphi} over all 12 neighbors and summing, the first-order
terms cancel by the symmetry of the FCC lattice. Computing the
second-order tensor over the 12 neighbor vectors \ensuremath{r_{j}:}

\begin{equation*}
\Sigma_{j} r_{j}\mu r_{j}\nu = 4\ell^{2}_{p} \delta \mu \nu
\end{equation*}

The discrete Laplacian therefore converges as:

\begin{equation*}
\Delta_{\mathrm{FCC}} \varphi \to 2\ell^{2}_{p} \nabla ^{2}\varphi
\end{equation*}

Substituting into the discrete wave equation and converting to physical
time t = \ensuremath{\tau} \ensuremath{\cdot} \ensuremath{\tau_{\mathrm{flop}}:}

\begin{equation*}
\partial ^{2}\varphi / \partial t^{2} = ( c^{2}_{\mathrm{TMS}} \cdot 2\ell^{2}_{p} / \tau^{2}_{\mathrm{flop}} ) \nabla ^{2}\varphi
\end{equation*}

Requiring this to equal \ensuremath{c^{2}\nabla ^{2}\varphi ,} and substituting \ensuremath{\tau_{\mathrm{flop}}} = \ensuremath{\ell_{p}/c} from
Section 3.3:

\begin{equation*}
c_{\mathrm{TMS}} = 1/\sqrt{2}
\end{equation*}

The factor \ensuremath{1/\sqrt{2}} = \ensuremath{\text{cos}(45^{\circ})} is the direct geometric projection factor of
FCC face-diagonal connectivity. The propagation speed is not a parameter
--- it is the cosine of the lattice's canonical angle. No free parameter
remains.

\end{proof}

\begin{corollary}[Consistency]
Substituting \ensuremath{c_{\mathrm{TMS}}} = \ensuremath{1/\sqrt{2}} into the physical wave
speed relation confirms:
\end{corollary}

\begin{equation*}
c_{\mathrm{physical}} = c_{\mathrm{TMS}} \cdot \sqrt{2} \cdot \ell_{p} / \tau_{\mathrm{flop}} = (1/\sqrt{2}) \cdot \sqrt{2} \cdot c = c
\end{equation*}

\subsection*{6.6 The Continuum Limit}

Substituting \ensuremath{c_{\mathrm{TMS}}} = \ensuremath{1/\sqrt{2}} and the Laplacian result into the discrete wave
equation:

\begin{equation*}
\partial ^{2}\varphi / \partial \tau^{2} = (1/2) \cdot 2\ell^{2}_{p} \nabla ^{2}\varphi = \ell^{2}_{p} \nabla ^{2}\varphi
\end{equation*}

Converting to physical time t = \ensuremath{\tau} \ensuremath{\cdot} \ensuremath{t_{p}} and using \ensuremath{\ell_{p}} = c \ensuremath{\cdot} \ensuremath{t_{p}:}

\begin{equation*}
\partial ^{2}\varphi / \partial t^{2} = c^{2} \nabla ^{2}\varphi
\end{equation*}

\begin{remark}[Stochasticity and the Continuum Limit]
\ensuremath{U_{\mathrm{sh}}} operates at every
flop as both the source of latent bits available for confirmation and
the stochastic disruptor of previously confirmed states. These dual
roles are not in tension --- they are symmetric expressions of the same
maximum entropy property established in Axiom 0. \ensuremath{U_{\mathrm{sh}}} has no attractors
and no preferred targets. Every bit in every region has equal
probability of disruption at every flop. No geometric bias exists in the
disruption mechanism itself.
\end{remark}

At the substrate level, \ensuremath{U_{\mathrm{sh}}} contributes two statistically symmetric
pressures to the confirmation field: a sourcing pressure drawing new
bits toward actualization, and a disrupting pressure returning confirmed
bits to the unactualized field. Because both pressures are drawn from
the same maximum entropy distribution with no preferred scale or
direction, their net contribution to the large-scale confirmation field
averages to zero across sufficiently many flops.

What survives to the continuum level is therefore not a noisy wave
equation with a stochastic forcing term --- it is the structured
confirmation cascade, which is the only process with geometric coherence
across scales. The wave equation describes this coherent structure. The
stochastic substrate does not add a forcing term to the continuum limit
because its two roles cancel symmetrically. The continuum field \ensuremath{\varphi} is the
large-scale average of a dynamically stable cascade instance. This
observable universe represents one such stable average over flops ---
not the unique or guaranteed outcome, but the instance whose causal
stack persists long enough to be inhabited.

\subsection*{6.7 The Wave Equation}

Recognizing flop time \ensuremath{\tau} as the emergent time coordinate t:

\begin{equation*}
\partial ^{2}\varphi / \partial t^{2} - c^{2}\nabla ^{2}\varphi = 0
\end{equation*}

This is the standard wave equation. It is not imposed on the TMS --- it
emerges necessarily from the discrete dynamics of a fault-tolerant FCC
confirmation lattice in the continuum limit. Physical wave phenomena are
the large-scale behavior of the confirmation field propagating through
the agent substrate.

\subsection*{6.8 The Polarization Field as Parallel Object}

\begin{remark}[The Polarization Field]
The confirmation density \ensuremath{\varphi} derived
above is one of two independent informational fields supported by the
FCC lattice. The polarization field \ensuremath{\psi} \ensuremath{\equiv} \ensuremath{\langle \chi \rangle} --- the local coarse-grained
mean of confirmed bit values, ranging over \ensuremath{[-1,} +1] --- propagates
by an analogous discrete wave equation on the same FCC lattice:
\end{remark}

\begin{equation*}
\partial ^{2}\psi / \partial \tau^{2} = c^{2}_{\mathrm{TMS}} \Delta_{\mathrm{FCC}} \psi
\end{equation*}

The propagation speed \ensuremath{c_{\mathrm{TMS}}} = \ensuremath{1/\sqrt{2}} is inherited from the lattice
geometry; the proof of Section 6.5 applies unchanged because both fields
propagate across the same nearest-neighbor structure. Both fields share
the same fault-tolerant FCC substrate. The two fields evolve
independently in the bulk and couple at confirmation events, where the
local incoming polarization wave amplitude biases the resolution of new
confirmations toward \ensuremath{\bar{B}} = 1 or \ensuremath{\bar{B}} = 0 against the unbiased stochastic kick
from \ensuremath{U_{\mathrm{sh}}.}

The dynamics of \ensuremath{\psi ,} its coupling to \ensuremath{\varphi} at confirmation events, the
emergence of computation from polarization-modulated propagation, the
structural correspondence between this coupling and the Born rule, and
the formal selection criterion for stable subsystem programs are
developed in Volume 1, Paper 2.

\section*{VII. Predictions and Directions for Subsequent Work}
\addcontentsline{toc}{section}{VII. Predictions and Directions for Subsequent Work}

\subsection*{7.1 The Generalized Uncertainty Principle: A Geometric Grounding}

The Generalized Uncertainty Principle (GUP) extends the standard
Heisenberg relation by introducing a minimum measurable length:

\begin{equation*}
\Delta x \cdot \Delta p \geq \hbar /2 \cdot (1 + \beta (\Delta p)^{2}/m^{2}_{p}c^{2})
\end{equation*}

where \ensuremath{\beta} is a dimensionless deformation parameter characterizing the
scale at which quantum gravitational effects become significant. The GUP
has been independently motivated by string theory, loop quantum gravity,
and quantum geometry (Hossenfelder, 2013). In all existing formulations,
\ensuremath{\beta} is a free parameter --- inserted phenomenologically rather than
derived from geometric first principles.

The TMS provides a geometric grounding for the minimum length that the
GUP encodes. The identification \ensuremath{\ell_{0}} \ensuremath{\equiv} \ensuremath{\ell_{p}} established in Section III has a
precise geometric meaning: below the agent scale, the emergent 3D
geometry does not exist. This is not a statement about measurement
difficulty or quantum fluctuations --- it is a statement about substrate
architecture. The floor on position resolution is not a statistical
lower bound but a geometric boundary. At \ensuremath{\ell_{p},} the continuous confirmation
field \ensuremath{\varphi} approaches its binary alphabet limit --- \ensuremath{\bar{B}} \ensuremath{\in} \{0, 1\} --- the
irreducible primitive of the TMS ontology. Spatial fuzziness at the
Planck scale is not noise added to geometry. It is the geometry of the
substrate itself becoming directly visible.

Examining a bit at the Planck scale means encountering not merely its
current confirmation state, but the full geometric and polarization
record of how it arrived there --- the layered causal history encoded in
the down-vector structure. What standard quantum mechanics describes as
wavefunction complexity is, in TMS terms, the geometric depth and
polarization complexity of the confirmation record. The uncertainty is
not epistemic. It is the irreducible causal narrative of the bit made
visible at the substrate boundary.

This geometric grounding is consistent with recent independent work
deriving GUP-type relations from spacetime microgeometry (Giné, 2026),
which proposes that generalized uncertainty principles may emerge as
effective kinematical relations induced by microscopic spacetime
geometry. The TMS provides an independent derivation from a more
fundamental starting point --- geometric primitives prior to spacetime,
rather than horizon microstructure within it.

\subsection*{7.2 A Novel Prediction: Causal-Depth-Dependent Uncertainty}

The TMS architecture implies a specific and novel extension of the
standard GUP that, to our knowledge, has no analog in existing
formulations.

In the TMS, every confirmed bit exists at the terminus of a causal chain
encoded in the down-vector structure. The confirmation field \ensuremath{\varphi} at any
point in the lattice carries not only the current confirmation state but
the full geometric record of how that state was reached --- the layered
history of prior confirmation events that produced the current position
in the causal stack. With the binary alphabet established in Section
1.4, each layer of that history additionally carries a specific
polarization value. We term the depth of this record the causal depth d
of the bit, and the specific pattern of polarizations across that depth
the polarization complexity of its history.

We conjecture that the effective deformation parameter \ensuremath{\beta} is not a
universal constant but a function of both:

\begin{equation*}
\beta = f(d, K_{\psi})
\end{equation*}

where d is the causal depth and \ensuremath{K_{\psi}} is the algorithmic complexity of
the polarization landscape encoding the bit's confirmation history. f is
monotonically increasing in both arguments. The precise form of f is a
prediction of the FCC confirmation geometry and is left for derivation
in subsequent work.

The physical interpretation is direct. A recently confirmed bit ---
shallow causal history, few prior layers, simple polarization record ---
has a substrate footprint close to the bare agent scale \ensuremath{\ell_{p}.} Its position
uncertainty is close to the geometric minimum. A bit embedded in a deep,
richly polarized causal structure has a larger substrate footprint ---
more of the lattice geometry and binary content is geometrically present
in its confirmation record. Its effective position uncertainty is
correspondingly larger than the bare minimum.

In standard quantum mechanical language: wavefunction complexity scales
with both causal depth and polarization complexity, and minimum position
uncertainty scales with wavefunction complexity. The TMS provides a
geometric derivation of this relationship from the down-vector structure
of the FCC causal stack and the binary alphabet recorded within it.

This prediction is testable in principle. If \ensuremath{\beta} scales with the joint
causal-depth and polarization-complexity profile, then quantum systems
with richer interaction histories --- more degrees of freedom, deeper
entanglement structure, and more complex binary records --- should
exhibit measurably larger GUP corrections than simpler systems at
equivalent energy scales. Recent work on GUP phenomenology in
macroscopic quantum systems and gravitational wave detectors
(Hossenfelder, 2013) provides experimental frameworks within which
causal-depth and polarization-complexity dependence could in principle
be distinguished from a universal \ensuremath{\beta .}

\subsection*{7.2a The Computational Irreducibility Connection}

The causal-depth dependence of \ensuremath{\beta} admits a deeper interpretation through
the lens of computational irreducibility (Wolfram, 2002). A
computationally irreducible system is one for which no shortcut exists
to determine its future state --- the only way to know where the system
is is to trace every step of its history.

In TMS terms, a bit with deep causal history and complex polarization
record is computationally irreducible in precisely this sense. The
down-vector chain encoding its confirmation record cannot be compressed
--- there is no shorter description of its current state than the full
sequence of prior confirmation events, including the specific binary
values resolved at each layer. This is not a limitation of our
knowledge. It is a property of the causal record itself.

The connection to Axiom 0 is direct. \ensuremath{U_{\mathrm{sh}}} is defined as algorithmically
incompressible with Kolmogorov complexity K(n) \ensuremath{\approx} n --- the maximum
entropy ground state from which all structure is actualized.
Confirmation events produce locally reducible structure by carving order
from that incompressibility. But the record of how that order was
produced --- the down-vector chain with its specific polarization
sequence --- inherits the incompressibility of its source. Deep causal
stacks with complex polarization landscapes are not merely geometrically
large. They are computationally irreducible in the strict Kolmogorov
sense.

The effective position uncertainty \ensuremath{\Delta x} \ensuremath{\geq} \ensuremath{\ell_{p}} \ensuremath{\cdot} f(d, \ensuremath{K_{\psi})} therefore has a
computational interpretation: measuring the position of a deeply
embedded bit to better than f(d, \ensuremath{K_{\psi})\cdot \ell_{p}} would require compressing an
irreducible computation. This is impossible by definition --- not in
principle, but as a consequence of the same Kolmogorov incompressibility
that characterizes \ensuremath{U_{\mathrm{sh}}} itself. The GUP minimum length in TMS is not
merely a geometric floor. It is the Planck-scale expression of
computational irreducibility applied to position measurement.

This connects the paper's opening axiom --- the incompressibility of \ensuremath{U_{\mathrm{sh}}}
--- directly to its closing prediction. The arc from Axiom 0 to the GUP
is a single unbroken chain:

\emph{\ensuremath{U_{\mathrm{sh}}} [K(n) \ensuremath{\approx} n] \ensuremath{\to} confirmation events \ensuremath{\to} binary alphabet \ensuremath{\to}
polarization-rich causal stack \ensuremath{\to} \ensuremath{\beta} = f(d, \ensuremath{K_{\psi})} \ensuremath{\to} GUP}

\subsection*{7.3 A Unified Geometric Account of Decoherence}

The same causal depth structure that predicts \ensuremath{\beta -\text{scaling}} offers a
geometric account of quantum decoherence. In standard quantum mechanics,
decoherence time depends on the number of environmental degrees of
freedom --- macroscopic objects decohere almost instantaneously because
they interact with enormously many environmental modes. The mechanism is
understood phenomenologically but lacks a geometric derivation from
first principles.

In the TMS, the causal stack provides a natural geometric analog. A bit
with shallow causal depth has a small substrate footprint --- it
interacts with few prior confirmation layers, maintains a simple touch
vector structure, and is relatively isolated from the full lattice
history. Such a bit maintains quantum coherence longer because its
geometric entanglement with the causal record is limited.

A bit with deep causal history has an extensive substrate footprint ---
its confirmation record spans many layers, its touch vector structure is
richly connected to the broader lattice, and it is geometrically
embedded in the full causal record. Such a bit decoheres rapidly because
its geometric connection to the lattice history is extensive --- there
are many more channels through which \ensuremath{U_{\mathrm{sh}}} disruption can reach and
collapse its coherent state.

Decoherence time, in this picture, is inversely proportional to causal
depth. Macroscopic objects decohere almost instantaneously because they
are geometrically ancient --- their causal stack is deep beyond measure.
Fundamental particles maintain quantum coherence because they are
geometrically shallow --- their confirmation records are brief relative
to the full causal stack depth.

This constitutes a unified geometric account of both the GUP minimum
length and quantum decoherence as consequences of a single structural
property --- causal depth and polarization complexity in the FCC
down-vector architecture. The derivation of the explicit functional
forms relating \ensuremath{\beta} and decoherence time to the joint causal-depth and
polarization-complexity profile is reserved for subsequent work.

\section*{Conclusion: A Complete Generative Mechanism}
\addcontentsline{toc}{section}{Conclusion: A Complete Generative Mechanism}

The Triadic Measurement Substrate presents a minimal, self-consistent
generative mechanism for wave dynamics and binary computational content
in three dimensions. Beginning from five axioms and the Principle of
Generative Economy, the following results emerge necessarily:

\begin{itemize}
\item
  Axioms 0--2 \ensuremath{\to} sphere primitive \ensuremath{\to} hexagonal 2D packing \ensuremath{\to} triadic
  confirmation
\item
  Triadic confirmation \ensuremath{\to} the \ensuremath{1\to 4} multiplier, exact and remainder-free by
  geometric necessity
\item
  Triadic confirmation pulling from the maximum-entropy substrate \ensuremath{\to}
  binary alphabet \ensuremath{\bar{B}} \ensuremath{\in} \{0, 1\} as the irreducible computational content
  of every confirmed bit, with polarization \ensuremath{\chi} \ensuremath{\in} \{\ensuremath{-1,} +1\} transported
  by every touch vector signal
\item
  Axiom 4 (causal anchoring) \ensuremath{\to} Z-displacement orthogonal by
  sphere-contact geometry \ensuremath{\to} ABCABC stacking \ensuremath{\to} FCC lattice
\item
  FCC lattice \ensuremath{\to} tetrahedral confirmation units \ensuremath{\to} [[4,2,2]]
  structural correspondence \ensuremath{\to} faithful encoding of value (LQ1) and
  causal-history (LQ2) \ensuremath{\to} fault-tolerant propagation of both
\item
  2D agent manifold \ensuremath{\to} minimum resolvable length identified with \ensuremath{\ell_{p}}
  (Planck length)
\item
  Planck length + speed of light \ensuremath{\to} \ensuremath{\tau_{\mathrm{flop}}} = \ensuremath{t_{p}} (Planck time)
\item
  FCC face-diagonal geometry \ensuremath{\to} \ensuremath{c_{\mathrm{TMS}}} = \ensuremath{1/\sqrt{2},} forced by lattice structure
\item
  \ensuremath{1\to 4} multiplier \ensuremath{\to} self-amplifying cascade regulated by \ensuremath{U_{\mathrm{sh}}} \ensuremath{\to}
  co-generated conjugate pair \ensuremath{(\varphi ,} \ensuremath{\Pi )} \ensuremath{\to} second-order dynamics
\item
  Fault-tolerant FCC lattice + second-order dynamics \ensuremath{\to} continuum limit \ensuremath{\to}
  \ensuremath{\partial ^{2}\varphi /\partial t^{2}} \ensuremath{-} \ensuremath{c^{2}\nabla ^{2}\varphi} = 0
\item
  Polarization field \ensuremath{\psi} as parallel object with identical propagation \ensuremath{\to}
  substrate for computational dynamics developed in Volume 1, Paper 2
\item
  Causal depth d and polarization complexity \ensuremath{K_{\psi}} in down-vector
  structure \ensuremath{\to} \ensuremath{\beta} = f(d, \ensuremath{K_{\psi})} \ensuremath{\to} GUP and decoherence from a single
  geometric property
\end{itemize}

Every quantity appearing in the TMS --- the absolute length scale, the
absolute time scale, the propagation speed, the binary alphabet --- is
fixed by the axioms and the Principle of Generative Economy without
external input. The confirmation rules and lattice connectivity that
govern the system's behavior are not assumptions imported from outside
the framework; they are the geometric consequences of the primitives
established in Section I. The model has no numerical degrees of freedom
left to tune.

The single fundamental rate of the system --- the flop --- governs
confirmation, synthesis, and propagation simultaneously. Flop time \ensuremath{\tau} is
not a background parameter but the counted succession of confirmation
events. Space is not a container but the accumulated causal record of
prior states. The [[4,2,2]] structural correspondence ensures
the system simultaneously maintains the present state, the binary value
of every confirmed bit, and the integrity of the causal record in every
flop. The wave equation emerges from the geometry of that record
propagating forward.

The TMS establishes that a discrete, informationally closed confirmation
substrate built from geometric primitives is both necessary and
sufficient to generate wave dynamics and a binary computational alphabet
from which physical laws may be derived. The paper opens with the
incompressibility of \ensuremath{U_{\mathrm{sh}}} (Axiom 0) and closes with a falsifiable
prediction --- causal-depth and polarization-complexity-dependent GUP
deformation --- grounded in that same incompressibility. The arc from
maximum entropy ground state to Planck-scale uncertainty is a single
unbroken geometric chain.

The derivation of further physical laws (electromagnetism, gravity), the
dynamics of the polarization field \ensuremath{\psi} and its coupling to confirmation
events, the explicit functional form of f(d, \ensuremath{K_{\psi}),} and the scaling
properties of the triadic confirmation rule across complex systems are
left for subsequent work in Volume 1.

\section*{Acknowledgments}
\addcontentsline{toc}{section}{Acknowledgments}

The development of this volume was substantially aided by extended
collaboration with several large language models used as critical
sounding boards, pedagogical resources, formal-rigor checkers, and
external working memory across many sessions during 2024--2026.

\textbf{Gemini (Google DeepMind)} was the long-running primary
collaborator across the framework's conceptual development, providing
physics pedagogy, structural critique, and ongoing review of the ideas
that became Volume 1. The author's path into the technical literature,
and the working-out of the framework's core conceptual moves over a long
period preceding formalization, were carried by extended dialogue with
Gemini.

\textbf{Claude (Anthropic)} was the primary collaborator across the
formalization, derivation tightening, simulation work, citation
auditing, and editorial passes that brought the volume to its present
form, including the sublattice bifurcation analysis empirically verified
in Paper 5 Appendix E.

\textbf{ChatGPT (OpenAI)}, \textbf{DeepSeek}, \textbf{Grok (xAI)}, and
\textbf{Granite (IBM)} contributed additional structural feedback and
review at various points.

All theoretical commitments, the choice of axioms, the Principle of
Generative Economy, the sublattice bifurcation insight, and the
framework's overall architecture are the author's; the AI systems' role
was to teach, formalize, critique, verify, and document.

\section*{References}
\addcontentsline{toc}{section}{References}
\begin{reflist}
Aitchison, C. T., \& Béri, B. (2025). Spacetime Spins: Statistical
mechanics for error correction with stabilizer circuits. arXiv preprint.
\url{https://arxiv.org/abs/2512.21991}

El Morsalani, M. (2025). Encoding Spacetime: A Comprehensive Review of
Holographic Quantum Codes. Preprint.

Fejes Tóth, L. (1940). Über einen geometrischen Satz. Mathematische
Zeitschrift, 46(1), 83--85. doi:10.1007/BF01181430.

Giné, J. (2026). On the Possibility of Quantum Gravity Emerging from
Geometry. arXiv preprint. \url{https://arxiv.org/abs/2602.16219}

Goldstein, H., Poole, C., \& Safko, J. (2002). Classical Mechanics (3rd
ed.). Addison-Wesley.

Gottesman, D. (1997). Stabilizer Codes and Quantum Error Correction.
Ph.D.~Thesis, California Institute of Technology.
\url{https://arxiv.org/abs/quant-ph/9705052}

Hamilton, W. R. (1834). On a General Method in Dynamics. Philosophical
Transactions of the Royal Society, 124, 247--308.

Hatcher, A. (2002). Algebraic Topology. Cambridge University Press.

Hossenfelder, S. (2013). Minimal Length Scale Scenarios for Quantum
Gravity. Living Reviews in Relativity, 16(1), 2.
doi:10.12942/lrr-2013-2.

Jaynes, E. T. (1957). Information Theory and Statistical Mechanics.
Physical Review, 106(4), 620--630. doi:10.1103/PhysRev.106.620.

Kauffman, S. (1993). The Origins of Order. Oxford University Press.

Kolmogorov, A. N. (1965). Three approaches to the quantitative
definition of information. Problems of Information Transmission, 1(1),
1--7.

Levin, M. (2019). The Computational Boundary of a Self: Developmental
Bioelectricity Drives Multicellularity and Scale-Free Cognition.
Frontiers in Psychology, 10, 2688. doi:10.3389/fpsyg.2019.02688.

Ma, Y., Tsao, D., \& Shum, H.-Y. (2022). On the Principles of Parsimony
and Self-Consistency for the Emergence of Intelligence. Frontiers of
Information Technology \& Electronic Engineering, 23, 1298--1323.
\url{https://arxiv.org/abs/2207.04630}

Pastawski, F., Yoshida, B., Harlow, D., \& Preskill, J. (2015).
Holographic quantum error-correcting codes: toy models for the
bulk/boundary correspondence. Journal of High Energy Physics, 2015(6),
149. doi:10.1007/JHEP06(2015)149. \url{https://arxiv.org/abs/1503.06237}

Rissanen, J. (1978). Modeling by shortest data description. Automatica,
14(5), 465--471. doi:10.1016/0005-1098(78)90005-5.

Shannon, C. E. (1948). A mathematical theory of communication. Bell
System Technical Journal, 27(3), 379--423.
doi:10.1002/j.1538-7305.1948.tb01338.x.

Thue, A. (1890). Om nogle geometrisk-taltheoretiske Theoremer.
Forhandlingerne ved de Skandinaviske Naturforskeres, 14, 352--353.

Wheeler, J. A. (1990). Information, physics, quantum: The search for
links. In W. H. Zurek (Ed.), Complexity, Entropy and the Physics of
Information (pp.~3--28). Addison-Wesley.

Wolfram, S. (2002). A New Kind of Science. Wolfram Media.
\end{reflist}

\clearpage


\end{document}
